\documentclass[11pt]{article}
\usepackage[margin=1in]{geometry}
\usepackage{amsmath,amssymb,amsthm,mathtools,booktabs,hyperref}
\hypersetup{hidelinks}
\newtheorem{theorem}{Theorem}
\newtheorem{lemma}[theorem]{Lemma}
\newtheorem{proposition}[theorem]{Proposition}
\newtheorem{corollary}[theorem]{Corollary}
\theoremstyle{remark}\newtheorem{remark}[theorem]{Remark}
\newcommand{\Z}{\mathbb Z}\newcommand{\N}{\mathbb N}
\newcommand{\R}{\mathbb R}\newcommand{\C}{\mathbb C}
\newcommand{\dd}{\,d}\newcommand{\comb}[1]{\delta_{#1}}
\title{Finite Euler reflection rigidity and a positive Poisson boundary\\
\large Paper C: Theorem Forge II}
\author{Observational SUCC research dossier}
\date{19 October 2026}
\begin{document}\maketitle
\begin{abstract}
A finite quotient of Euler factors with real bases greater than one cannot
satisfy reflection about a nonzero real center unless its two base multisets
are identical. Equal cardinality is a conclusion, not an assumption. We give
an exact reflection failure for a density-one Beurling counterfeit. For infinite
sources we prove a separate sufficient theorem: a positive tempered measure
with a normalized atom at zero, a unit gap and Fourier self-duality is the
integer Dirac comb. An explicit contour-decay hypothesis converts the usual
Gaussian-completed functional equation into this rigidity criterion, including
for generalized real exponents. No growth estimate is inferred from a formal
functional equation. These are source-identification results, not a proof of
the Riemann hypothesis. Scholarly novelty remains unverified.
\end{abstract}

\section{Contract and source boundary}
The finite theorem concerns finite multisets of real numbers greater than
one, with multiplicities retained. The infinite theorem concerns positive
source measures and states all analytic hypotheses separately. The earlier
repository result ASM-01 constructs, beyond every finite observation horizon,
a collision-free Beurling source with density one, the ordinary prime number
theorem, residue one and the same critical-strip zeta divisor as ordinary
arithmetic, but a failed additive-rank braid. Its finite Euler quotient is
holomorphic and nonzero on the right half-plane; it does not preserve
the completed functional equation. We identify a precise extra condition that
excludes this family and a sufficient infinite-source extension.

The base is PR179, commit
\texttt{81684b4b8956f9f5da56c376ba7e4339277eb665}; the predecessor proof is
\texttt{research/2026-10-18/ARITHMETIC\_SOURCE\_MINIMALITY.md}.
Nothing below assumes that sibling PR178 has been merged. No zero ordinate
is an input to a construction or a test.

\section{Finite factors: no cardinality assumption}
For finite multisets $U,V\subset(1,\infty)$ define the meromorphic function
\[
 F(s)=\frac{\prod_{u\in U}(1-u^{-s})}{\prod_{v\in V}(1-v^{-s})},
 \qquad u^{-s}=e^{-s\log u}.
\]
Empty products equal one. Factors and their multiplicities are literal;
coincident powers of different bases are not identified.

\begin{theorem}[Finite reflection rigidity]\label{thm:finite}
Let $c\in\R\setminus\{0\}$. If $F(s)=F(c-s)$ as meromorphic functions on
$\C$, then $F\equiv1$ and $U=V$ as multisets. In particular the statement
holds for $c=1$, without assuming $|U|=|V|$.
\end{theorem}
\begin{proof}
Put $\epsilon=(-1)^{|U|+|V|}$ and
$C=\prod_{u\in U}u/\prod_{v\in V}v>0$. The elementary factor identity
$1-u^x=-u^x(1-u^{-x})$ gives, for real $t$ sufficiently large,
\begin{equation}\label{eq:inversion}
 F(c-t)=\epsilon C^{t-c}F(t-c).
\end{equation}
Both $F(t)$ and $F(t-c)$ are strictly positive and tend to one as
$t\to\infty$. Reflection and \eqref{eq:inversion} first imply
$\epsilon=1$, and then
$C^{t-c}=F(t)/F(t-c)\to1$, so $C=1$. Consequently $F(t)=F(t-c)$ for
all sufficiently large $t$. Whether $c$ is positive or negative, this gives
$F(t)=F(t+n|c|)$ for every $n\geq0$ and every sufficiently large $t$.
Taking $n\to\infty$ proves $F(t)=1$. The meromorphic identity theorem
gives $F\equiv1$.

Cancel equal bases on the two sides with their full multiplicities. If any
base remains, let $b$ be the least one and let $m_b\ne0$ be its signed
multiplicity (numerator minus denominator). On the positive real axis the
real logarithm is defined and
\[
 \log F(t)=\sum_a m_a\log(1-a^{-t}),\qquad
 \lim_{t\to\infty}b^t\log F(t)=-m_b.
\]
Indeed $\log(1-z)=-z+O(z^2)$ as $z\to0$; every remaining base $a>b$
contributes a quantity tending to zero after multiplication by $b^t$, and
the sum is finite. This contradicts $F\equiv1$. Thus no base remains.
\end{proof}

There is also a direct cardinality check when $c=1$. Each factor has a
simple zero at zero, so $\operatorname{ord}_0F=|U|-|V|$. Since $F(1)$ is
finite and strictly positive, reflection forces this order to vanish.
This independently handles unequal counts. Possible poles on the imaginary
axis cause no defect in the proof: its real-axis ratios have no zero or pole,
and all other locations are handled by meromorphic continuation.

\begin{remark}[Sharp excluded center]
For $U=\{2,6\}$ and $V=\{3,4\}$, the counts and products agree and
\eqref{eq:inversion} gives $F(s)=F(-s)$. The quotient is nonconstant by the
least-base argument. Thus the restriction $c\ne0$ is necessary. Equal
counts and equal products alone are insufficient at $c=1$.
\end{remark}

\subsection{An exact counterfeit rejected by completion}
Delete ordinary primes $p=3,q=11$ and introduce primitive values
$a=9/2,b=77/17$. Retain every other ordinary prime and count formal
commutative words. The rational valuations at the deleted primes are
\[
 (v_3(a),v_{11}(a))=(2,0),\qquad
 (v_3(b),v_{11}(b))=(0,1),
\]
so distinct words have distinct values. For $\Re s>1$ its zeta function is
$Z_B(s)=\zeta(s)F(s)$ with $U=\{3,11\}$ and $V=\{9/2,77/17\}$.
The identity
\[
 (1-a^{-1})(1-b^{-1})=(1-3^{-1})(1-11^{-1})
\]
gives $F(1)=1$. None of these factors vanishes in $\Re s>0$, so the same
critical-strip divisor and residue one follow from ordinary zeta there.
But exact rational evaluation yields
\begin{equation}\label{eq:counterfeit}
 F(2)=\frac{504}{517},\qquad F(-1)=\frac{34}{21}.
\end{equation}
In particular the same Gaussian completion cannot obey reflection. More
generally, if $\Lambda_B=\Lambda_\zeta F$ and both completed functions
obey reflection, meromorphic cancellation of the nonzero function
$\Lambda_\zeta$ forces $F(s)=F(1-s)$. Theorem~\ref{thm:finite} then
forbids every nontrivial finite alteration of primitive bases, including
alterations with repeated generators or coincident powers.

\section{A positive Poisson rigidity lemma}
Use the Fourier convention
$\widehat f(\xi)=\int_\R f(x)e^{-2\pi i x\xi}\dd x$ and its
extension to tempered distributions. Write
$\comb{a\Z}=\sum_{k\in\Z}\delta_{ak}$.

\begin{theorem}[Unit gap and positive self-duality]\label{thm:gap}
Let $\mu$ be a positive tempered Radon measure on $\R$ with
\[
 \mu(\{0\})=1,\qquad \mu((-1,1)\setminus\{0\})=0,
 \qquad \widehat\mu=\mu.
\]
Then $\mu=\comb{\Z}$. Equivalently, an arbitrary positive mass $\rho$
at zero gives $\mu=\rho\comb{\Z}$ under the other two hypotheses.
\end{theorem}
\begin{proof}
Set
\[
 \phi(x)=\left(\frac{\sin\pi x}{\pi x}\right)^2,\quad \phi(0)=1,
 \qquad T(\xi)=(1-|\xi|)_+.
\]
The Fourier identity $\widehat\phi=T$ follows by squaring the inverse
transform of $\mathbf1_{[-1/2,1/2]}$ and convolving that interval with
itself. Since $\phi$ is not Schwartz, we do not pair it with an arbitrary
tempered measure before proving integrability. Instead put
\[
 \phi_\varepsilon(x)=e^{-\pi\varepsilon x^2}\phi(x),\qquad
 G_\varepsilon(\xi)=\varepsilon^{-1/2}e^{-\pi\xi^2/\varepsilon}.
\]
For $\varepsilon>0$ these are Schwartz functions and
$\widehat{\phi_\varepsilon}=T*G_\varepsilon\geq0$. Self-duality gives
\begin{equation}\label{eq:smoothing}
 \int\phi_\varepsilon\dd\mu
   =\int(T*G_\varepsilon)\dd\mu.
\end{equation}
The right side tends to $\int T\dd\mu=1$. Here are the needed tail
details. A positive tempered measure has polynomial mass growth: testing
against a nonnegative Schwartz function that is at least one on $[-1,1]$
and dilating it bounds $\mu([-R,R])$ by a fixed polynomial in $R\geq1$.
On compact sets $0\leq T*G_\varepsilon\leq1$ and pointwise convergence
to the continuous function $T$ holds, including at $\pm1$. For
$|\xi|\geq2$ and $0<\varepsilon\leq1$, integration over $[-1,1]$
and $|\xi-y|\geq|\xi|-1$ give
\[
 (T*G_\varepsilon)(\xi)
 \leq 2\varepsilon^{-1/2}
        e^{-\pi(|\xi|-1)^2/\varepsilon}
 \leq2e^{-\pi(|\xi|-1)^2}.
\]
The last inequality follows by maximizing
$v^{1/2}e^{-\pi d^2v}$ for $v\geq1,d\geq1$. This Gaussian envelope
is integrable against a measure of polynomial growth. Dominated convergence
therefore applies to the right side of \eqref{eq:smoothing}.

On the left, $\phi_\varepsilon\uparrow\phi$ as
$\varepsilon\downarrow0$. Monotone convergence proves
$\int\phi\dd\mu=1$, now including finiteness. The atom at zero
already contributes one. Positivity forces all other mass into the zeros of
$\phi$, namely the nonzero integers. Thus $\operatorname{supp}\mu\subset\Z$.
Because $\widehat\mu=\mu$ is a measure supported on $\Z$,
$(e^{2\pi i\xi}-1)\widehat\mu=0$. The Fourier translation rule then
implies that $\mu$ is invariant under translation by one. All integer atoms
therefore have the same mass, equal to one. Rescaling proves the final claim.
\end{proof}

\subsection{Decisive mutations and an exact local obstruction}
The gap cannot be omitted. The positive measure
\[
 \mu_* =\frac23\comb{2\Z}+\frac13\comb{\Z/2}
\]
has $\mu_*(\{0\})=1$ and is self-dual, because
$\widehat{\comb{a\Z}}=a^{-1}\comb{\Z/a}$. It has an atom of weight
$1/3$ at $1/2$, so it is not the integer comb. Its positive-frequency
Dirichlet series is
\[
 Z_*(s)=\frac{2^s+2^{1-s}}3\zeta(s).
\]
The multiplier is reflection invariant and equals one at $s=1$; this
control retains residue and completion. It is not a Beurling word source
with all primitive values greater than one, since it has frequencies below
one. Conversely, $\comb{\Z}+\delta_{3/2}+\delta_{-3/2}$ retains positivity,
normalization and the gap but fails self-duality by the proof above.

Suppose a positive unit-gap measure has atoms of mass at least $w>0$ at
$\pm x$, where $x\geq1$ is not an integer. Write
$d=\operatorname{dist}(x,\Z)\in(0,1/2]$. Concavity of sine on
$[0,\pi/2]$ gives $\sin(\pi d)\geq2d$, hence
\begin{equation}\label{eq:obstruction}
 \pi^2\left(\int\phi\dd\mu-\mu(\{0\})\right)
 \geq \frac{8wd^2}{x^2}>0,
\end{equation}
with an infinite left side allowed. This contradicts self-duality without
an infinite numerical sum. The counterfeit atom $x=9/2,w=1$ already gives
the exact pair contribution $8/81$ to the left side. The rational code
certifies this lower bound; it does not certify the Fourier identity by a
finite sample.

\section{What an infinite-source completion must actually supply}
Let $B$ be a locally finite multiset in $[1,\infty)$, with multiplicities
$m_b\in\N$ and $m_1=1$. A Beurling word multiset is a special case: take
primitive real bases greater than one, finitely many at most each bound,
and count all their finite commutative words, including the unit. Assume
$Z_B(s)=\sum_{b\in B}m_b b^{-s}$ converges absolutely in some right
half-plane, where the sum here and below is over distinct values.
No Euler product is required for the first conclusion of the next theorem.

\begin{theorem}[A sufficient infinite-source rigidity contract]
\label{thm:infinite}
Define $\Lambda_B(s)=\pi^{-s/2}\Gamma(s/2)Z_B(s)$ in its initial
half-plane. Suppose that for some $c>1$ in that half-plane:
\begin{enumerate}
\item $\Lambda_B$ extends meromorphically to $\C$, with exactly two
simple poles, at $0,1$, with residues $-\rho,\rho$, where $\rho>0$;
\item $\Lambda_B(s)=\Lambda_B(1-s)$;
\item uniformly for $1-c\leq\sigma\leq c$ as $|t|\to\infty$,
$|\Lambda_B(\sigma+it)|=O((1+|t|)^{-2})$.
\end{enumerate}
Then $\rho=1$ and $B$ is the positive integer multiset, each integer
having multiplicity one. If $B$ is a Beurling word multiset, its primitive
multiset is exactly the ordinary primes, with multiplicity one.
\end{theorem}
\begin{proof}
Absolute convergence at $c$ implies
$\sum_{b\leq X}m_b\leq X^c Z_B(c)$, so the symmetric source measure
below is tempered and all Gaussian sums converge. Put
$f(t)=\sum_bm_b e^{-\pi b^2t}$ for $t>0$. Mellin inversion of the
exponential, followed by absolutely justified summation, gives
\begin{equation}\label{eq:mellin}
 f(t)=\frac1{4\pi i}\int_{\Re s=c}\Lambda_B(s)t^{-s/2}\dd s.
\end{equation}
The absolute justification on this line follows from
$\sum_bm_b b^{-c}<\infty$ and the integrability of
$\Gamma((c+iy)/2)$ in $y$. Shift the contour to $\Re s=1-c$.
The uniform decay hypothesis makes the horizontal edges vanish and both
vertical integrals absolutely convergent. The residues at $1$ and $0$
give respectively $\rho t^{-1/2}/2$ and $-\rho/2$. Substituting
$s=1-w$ in the remaining vertical integral and using reflection yields
\[
 f(t)=t^{-1/2}f(1/t)+\frac\rho2(t^{-1/2}-1).
\]
Consequently $\theta(t)=\rho+2f(t)$ obeys
$\theta(t)=t^{-1/2}\theta(1/t)$.

Define the positive even tempered measure
\[
 \mu=\rho\delta_0+\sum_bm_b(\delta_b+\delta_{-b}).
\]
For $g_t(x)=e^{-\pi tx^2}$, the Gaussian transform formula and theta
reciprocity imply $\langle\mu-\widehat\mu,g_t\rangle=0$ for every
$t>0$. The even-distribution separation argument in
Lemma~\ref{lem:gaussians} below gives $\mu=\widehat\mu$.
Theorem~\ref{thm:gap}, after rescaling by $\rho$, gives
$\mu=\rho\comb{\Z}$. The source atom at $1$ has mass $m_1=1$,
so $\rho=1$. Thus $B=\N_{>0}$ with multiplicity one.

For a Beurling word source, each primitive must now be an integer at least
two. Each ordinary prime must occur as a primitive, since a product of
integers at least two can equal a prime only when it is a single factor.
A repeated primitive or any composite primitive would give at least two
different words at the same integer (using ordinary prime factors),
contradicting multiplicity one. This identifies the primitive multiset.
\end{proof}

\begin{lemma}[Centered Gaussians separate even tempered distributions]
\label{lem:gaussians}
If $T\in\mathcal S'(\R)$ is even and
$\langle T,e^{-\pi tx^2}\rangle=0$ for every $t>0$, then $T=0$.
\end{lemma}
\begin{proof}
The map $t\mapsto e^{-\pi tx^2}$ is smooth with values in Schwartz
space on $t>0$. Differentiating at $t=1$ shows that $T$ annihilates
$x^{2k}e^{-\pi x^2}$ for every $k\geq0$. It thus annihilates every
even dilated Hermite function with Gaussian factor $e^{-\pi x^2}$,
using the oscillator $-d^2/dx^2+4\pi^2x^2$. Its polynomial factor is even.
The corresponding dilated Hermite
expansions of Schwartz functions converge in the Schwartz topology, so
these functions span a dense subspace of the even Schwartz functions.
For completeness, the topology assertion follows by applying powers of the
harmonic oscillator: the Hermite coefficients of a Schwartz function
decay faster than every inverse power of the index, while every Schwartz
seminorm of a Hermite function grows at most polynomially, by the raising
and lowering identities. Therefore the expansion and its derivatives with
polynomial weights converge in every Schwartz seminorm. Continuity of $T$
shows that it annihilates all even Schwartz functions. Its evenness makes
it annihilate odd Schwartz functions as well, proving $T=0$.
\end{proof}

\paragraph{Analytic boundary.}
The strip decay in Theorem~\ref{thm:infinite} is a strong, sufficient
condition chosen to make every contour step explicit. It is not deduced
from the displayed functional equation, residue data, finite Euler jets or
finite observations. Absolute convergence on a right half-plane suffices
for polynomial source growth, so a separate density-one hypothesis is not
needed. The theorem is not a classification without growth conditions,
nor a claim that every generalized prime zeta has the ordinary completion.
It excludes all nonordinary positive Beurling word sources satisfying its
complete contract; it does not establish that contract from the repository's
observer axioms.

\section{Prior art, status and the next falsifier}
Hilberdink and Lapidus~\cite{HL} develop generalized prime zeta functions
and relate Mellin functional equations to modular identities under explicit
analytic hypotheses (Theorem 3.2). Their paired completions need not be the
same Gaussian completion. This supplies context, not a classification
premise for Theorem~\ref{thm:infinite}.

Burnol~\cite{Burnol} recalls Hamburger's converse for ordinary integer
Dirichlet series and gives extensions, including generalized exponents.
One cannot apply the ordinary-support theorem to reconstruct integer
support by assuming that support at the outset. Our positive-gap proof
makes that step explicit under the stronger contour contract stated here.

Baake, Spindeler and Strungaru~\cite{BSS} study Fourier eigenmeasures far
beyond a single lattice and record the lattice Poisson formula and Hermite
framework. Self-duality alone is therefore not our identifying premise.
The added positive unit gap yields an extremal sinc argument of classical
Fourier type; priority for this particular formulation is unverified.

All three theorems have complete proofs above and exact adverse controls.
Their status is \emph{proved locally, scholarly novelty unverified}. The
finite tests exercise arithmetic identities and witnesses, not universal
analytic claims. The cheapest falsifier for finite camouflage is
\eqref{eq:counterfeit}; for a noninteger positive source atom it is
\eqref{eq:obstruction}. The remaining source lemma is to derive the
full normalized Gaussian-completed continuation and strip control from
independent source axioms, or to exhibit a nonordinary source that passes
those axioms while violating one of these analytic hypotheses. Merely
renaming reflection as a source law does not close that dependency.

\begin{thebibliography}{9}
\bibitem{HL} T. Hilberdink and M. L. Lapidus,
\emph{Beurling Zeta Functions, Generalised Primes, and Fractal Membranes},
Acta Appl. Math. 94 (2006), 21--48.
\href{https://doi.org/10.1007/s10440-006-9063-0}{doi:10.1007/s10440-006-9063-0}.
Primary full text: \url{https://arxiv.org/abs/math/0410270}.
\bibitem{Burnol} J.-F. Burnol,
\emph{Deux extensions de Th\'eor\`emes de Hamburger}, 2011, version 2.
\url{https://arxiv.org/abs/1106.4749}.
\bibitem{BSS} M. Baake, T. Spindeler and N. Strungaru,
\emph{On Eigenmeasures Under Fourier Transform},
J. Fourier Anal. Appl. 29 (2023), article 65.
\href{https://doi.org/10.1007/s00041-023-10045-z}{doi:10.1007/s00041-023-10045-z}.
\end{thebibliography}
\end{document}
